% ---------------------------------------------------------------
% Adaptado de https://pt.overleaf.com/latex/templates/template-ufabc-dissertacao/zmwgdkcsrxjb
% por Francisco de Assis Zampirolli, 2023
% ---------------------------------------------------------------

% ---------------------------------------------------------------
% Modelo LaTex para dissertação e tese do programa de Pós Graduação em Ciência da Computação da UFABC
% ---------------------------------------------------------------

\documentclass[
	% -- opções da classe memoir --
	12pt,					% tamanho da fonte
	openright,				% capítulos começam em pág ímpar (insere página vazia caso preciso)
	oneside,					% para impressão em verso e anverso. Oposto a oneside
	a4paper,					% tamanho do papel. 
	% -- opções da classe abntex2 --
	%chapter=TITLE,			% títulos de capítulos convertidos em letras maiúsculas
	%section=TITLE,			% títulos de seções convertidos em letras maiúsculas
	%subsection=TITLE,		% títulos de subseções convertidos em letras maiúsculas
	%subsubsection=TITLE,	% títulos de subsubseções convertidos em letras maiúsculas
	% -- opções do pacote babel --
	english,					% idioma adicional para hifenização
	%french,					% idioma adicional para hifenização
	%spanish,				% idioma adicional para hifenização
	brazil					% o último idioma é o principal do documento
	]{abntex2}

% ---------------------
% Pacotes OBRIGATÓRIOS
% ---------------------
\usepackage{lmodern}				% Usa a fonte Latin Modern			
\usepackage[T1]{fontenc}			% Selecao de codigos de fonte.
\usepackage[utf8]{inputenc}		% Codificacao do documento (conversão automática dos acentos)
\usepackage{lastpage}			% Usado pela Ficha catalográfica
\usepackage{indentfirst}			% Indenta o primeiro parágrafo de cada seção.
\usepackage{color}				% Controle das cores

\usepackage{graphicx,graphicx}	% Inclusão de gráficos
\usepackage{epsfig,subfig}		% Inclusão de figuras
\usepackage{microtype} 			% Melhorias de justificação
% ---------------------

% Colors package
\usepackage[dvipsnames]{xcolor}
\usepackage {soul}
\setul{0.3ex}{0.2ex}
 
% ---------------------
% Pacotes ADICIONAIS
% ---------------------
\usepackage{lipsum}						% Geração de dummy text
\usepackage{amsmath,amssymb,mathrsfs}	% Comandos matemáticos avançados 
\usepackage{setspace}  					% Para permitir espaçamento simples, 1 1/2 e duplo
\usepackage{verbatim}					% Para poder usar o ambiente "comment"
\usepackage{tabularx} 					% Para poder ter tabelas com colunas de largura auto-ajustável
\usepackage{afterpage} 					% Para executar um comando depois do fim da página corrente
\usepackage{url} 						% Para formatar URLs (endereços da Web)
% ---------------------

\usepackage{quoting}

\usepackage{hyperref}

% ---------------------
% Pacotes de CITAÇÕES
% ---------------------
\usepackage[brazilian,hyperpageref]{backref}	% Paginas com as citações na bibl
%\usepackage[alf,abnt-etal-cite=4]{abntex2cite}				% Citações padrão ABNT (alfa)
\usepackage[alf,abnt-etal-cite=4,abnt-etal-list=0,abnt-etal-text=emph]{abntex2cite}
%\usepackage[num]{abntex2cite}				% Citações padrão ABNT (numericas)
%\usepackage[style=alphabetic,sorting=nyt,sortcites=true,maxbibnames=9,autopunct=true,babel=hyphen,hyperref=true,abbreviate=false]{biblatex}
% ---------------------


% Configurações de CITAÇÕES para abntex2
% --- 
% CONFIGURAÇÕES DE PACOTES
% --- 

% ---
% Configurações do pacote backref
% Usado sem a opção hyperpageref de backref
\renewcommand{\backrefpagesname}{Citado na(s) página(s):~}
% Texto padrão antes do número das páginas
\renewcommand{\backref}{}
% Define os textos da citação
\renewcommand*{\backrefalt}[4]{
	\ifcase #1 %
		Nenhuma citação no texto.%
	\or
		Citado na página #2.%
	\else
		Citado #1 vezes nas páginas #2.%
	\fi}%
% ---

% Inclusão de dados para CAPA e FOLHA DE ROSTO (título, autor, orientador, etc.)
\input{extras/dados}

% Inclui Configurações de aparência do PDF Final
%  Configurações de aparência do PDF final
% NÃO ALTERAR!!!

% alterando o aspecto da cor azul
\definecolor{blue}{RGB}{41,5,195}

% informações do PDF
\makeatletter
\hypersetup{
     	%pagebackref=true,
		pdftitle={\@title}, 
		pdfauthor={\@author},
    		pdfsubject={\imprimirpreambulo},
	    pdfcreator={LaTeX with abnTeX2},
		pdfkeywords={abnt}{latex}{abntex}{abntex2}{trabalho acadêmico}, 
		colorlinks=true,       		% false: boxed links; true: colored links
    		linkcolor=blue,          	% color of internal links
    		citecolor=blue,        		% color of links to bibliography
    		filecolor=magenta,      		% color of file links
		urlcolor=blue,
		bookmarksdepth=4
} 
\makeatother
% --- 

% O tamanho da identação do parágrafo é dado por:
\setlength{\parindent}{1.3cm}

% Controle do espaçamento entre um parágrafo e outro:
\setlength{\parskip}{0.2cm}  % tente também \onelineskip

% ---------------------
% Compila o índice
% ---------------------
\makeindex
% ---------------------

%%%%%%%%%%%%%%%%%%%%%%%%%%%
%%  INICIO DO DOCUMENTO  %%
%%%%%%%%%%%%%%%%%%%%%%%%%%%
\begin{document}

% Retira espaço extra obsoleto entre as frases.
\frenchspacing

% ----------------------------------------------------------
% ELEMENTOS PRÉ-TEXTUAIS (Capa, Resumo, Abstract, etc.)
% ----------------------------------------------------------
\pretextual

% Capa
% ---
% Impressão da Capa
% ---
  \begin{capa}%
    \begin{figure}[h!]%
        \centering%
        \includegraphics[scale=1.2]{figs/logo.png}%
      \end{figure}%
    \center
	\ABNTEXchapterfont\large{Universidade Federal do ABC \\ Centro de Matemática, Computação e Cognição\\ Área: Processamento de Imagens e Visão Computacional}
	%\vspace{1.5cm}

    \vfill
    \ABNTEXchapterfont\bfseries\LARGE\imprimirtitulo
    \vfill

	%\vfill
	\ABNTEXchapterfont\large\imprimirautor
	\vfill
%
	%Número de Ordem : XXXX
	
    \large\imprimirlocal, \large\imprimirdata

    \vspace*{1cm}
  \end{capa}
% ---

% % Folha de rosto (o * indica que haverá a ficha bibliográfica)
% \imprimirfolhaderosto*

% % Imprimir Ficha Catalografica
% % ---
% Ficha Catalográfica
% ---
% Isto é um exemplo de Ficha Catalográfica, ou ``Dados internacionais de
% catalogação-na-publicação''. Você pode utilizar este modelo como referência. 
% Porém, talvez a biblioteca lhe fornece um PDF
% com a ficha catalográfica definitiva após a defesa do trabalho. Quando estiver
% com o documento, salve-o como PDF no diretório do seu projeto e substitua todo
% o conteúdo de implementação deste arquivo pelo comando abaixo:
%
% \begin{fichacatalografica}
%     \includepdf{fig_ficha_catalografica.pdf}
% \end{fichacatalografica}
\begin{fichacatalografica}
	\vspace*{\fill}					% Posição vertical
	\hrule							% Linha horizontal
	\begin{center}					% Minipage Centralizado
	\begin{minipage}[c]{12.5cm}		% Largura
	
	\imprimirautor
	
	\hspace{0.5cm} \imprimirtitulo  / \imprimirautor. --
	\imprimirlocal, \imprimirdata-
	
	\hspace{0.5cm} \pageref{LastPage} p. : il. (algumas color.) ; 30 cm.\\
	
	\hspace{0.5cm} \imprimirorientadorRotulo~\imprimirorientador\\
	
	\hspace{0.5cm}
	\parbox[t]{\textwidth}{\imprimirtipotrabalho~--~\imprimirinstituicao,
	\imprimirdata.}\\
	
	\hspace{0.5cm}
		1. Palavra-chave1.
		2. Palavra-chave2.
		%I. Orientador.
		I. Universidade Federal do ABC.
		II. Centro de Matemática, Computação e Cognição.
		III. Professor Titular\\ 			
	
	\hspace{8.75cm} CDU 02:141:005.7\\
	
	\end{minipage}
	\end{center}
	\hrule
\end{fichacatalografica}
% ---

% % Inserir Folha de Aprovação
% % ---
% Assinaturas
% ---
% Este é um exemplo de folha de aprovação, elemento obrigatório da NBR
% 14724/2011 (seção 4.2.1.3). Você pode utilizar este modelo até a aprovação
% do trabalho. Após isso, substitua todo o conteúdo deste arquivo por uma
% imagem da página assinada pela banca com o comando abaixo:
%
% \includepdf{folhadeaprovacao_final.pdf}
%
\begin{folhadeaprovacao}

  \begin{center}
    {\ABNTEXchapterfont\large\imprimirautor}

    \vspace*{\fill}\vspace*{\fill}
    \begin{center}
      \ABNTEXchapterfont\bfseries\Large\imprimirtitulo
    \end{center}
    \vspace*{\fill}
    
    \hspace{.45\textwidth}
    \begin{minipage}{.5\textwidth}
        \imprimirpreambulo
    \end{minipage}%
    \vspace*{\fill}
   \end{center}
        
 % Isso na versao final do trabalho!!!       
   Trabalho aprovado. \imprimirlocal, XX de fevereiro de 2024:

   %\assinatura{\textbf{\imprimirorientador} \\ Orientador} 
   %\assinatura{\textbf{\imprimircoorientador} \\ Co-Orientador} 
   \assinatura{\textbf{Professor} \\ Presidente - Interno}      
   \assinatura{\textbf{Professor} \\ Convidado 1}
   \assinatura{\textbf{Professor} \\ Convidado 2}
   \assinatura{\textbf{Professor} \\ Convidado 3}
   \begin{center}
    \vspace*{0.5cm}
    {\large\imprimirlocal}
    \par
    {\large\imprimirdata}
    \vspace*{1cm}
  \end{center}
  
\end{folhadeaprovacao}
% ---

% % Dedicatória
% % ---
% Dedicatória
% ---
\begin{dedicatoria}
   \vspace*{\fill}
   \centering
   \noindent
   \textit{À minha família, Cristina, Rafael e Eduardo, com muito amor.} \vspace*{\fill}
\end{dedicatoria}
% ---

% % Agradecimentos
% % ---
% Agradecimentos
% ---
\begin{agradecimentos}


Agradeço aos meus colegas da UFABC, professores, técnicos e alunos, que foram fundamentais para chegar até nesse pedido de promoção para professor titular do magistério superior, título desejado a muito tempo.

Aos meus amigos ...

Aos professores ...

À XXXXXX pelo apoio financeiro para realização deste trabalho de pesquisa.

\end{agradecimentos}
%% ---

% % Epígrafe
% % ---
% Epígrafe
% ---
\begin{epigrafe}
    \vspace*{\fill}
	\begin{flushright}
		\textit{``Não sei o que, \\
		          não sei o que,\\
                  não sei o que lá.''\\
		          (Autor Desconhecido)}
	\end{flushright}
\end{epigrafe}
% ---

% Resumo e Abstract
% ---
% RESUMOS
% ---

% RESUMO em português
\setlength{\absparsep}{18pt} % ajusta o espaçamento dos parágrafos do resumo
\begin{resumo}
\input{texLattes/Resumo.tex} \\
 \textbf{Palavras-chaves}: latex. abntex. editoração de texto.
\end{resumo}

% % ABSTRACT in english
% \begin{resumo}[Abstract]
%  \begin{otherlanguage*}{english}
%    This is the english abstract.

%    \vspace{\onelineskip}
 
%    \noindent 
%    \textbf{Keywords}: latex. abntex. text editoration.
%  \end{otherlanguage*}
% \end{resumo}

% Lista de ilustrações
\pdfbookmark[0]{\listfigurename}{lof}
\listoffigures*
\cleardoublepage

% Lista de tabelas
\pdfbookmark[0]{\listtablename}{lot}
\listoftables*
\cleardoublepage

% Lista de abreviaturas e siglas
% \begin{siglas}
%   \item[FAPESP] Fundação de Amparo à Pesquisa do Estado de São Paulo
%   \item[FEEC] Faculdade de Engenharia Elétrica e de Computação
%   \item[CMCC] Centro de Matemática, Computação e Cognição
%   \item[CNPq] Conselho Nacional de Desenvolvimento Científico e Tecnológico
%   \item[CPCE] Comissão de Promoção à Classe E
%   \item[CPPD] Comissão Permanente de Pessoal Docente
%   \item[IME] Instituto de Matemática e Estatística
%   \item[INEP] Instituto Nacional de Estudos e Pesquisas Educacionais Anísio Teixeira
%   \item[MEC] Ministério da Educação
%   %\item[MMach] Mathematical Morphology toolbox for the KHOROS
%   %\item[SisGRAFO] Sistema Gráfico de Otimização para Suporte
% à Decisão
%   \item[UFABC] Universidade Federal do ABC
%   \item[UFES] Universidade Federal do Espírito Santo
%   \item[UNICAMP] Universidade Estadual de Campinas
%   \item[USP] Universidade de São Paulo
% \end{siglas}

% % Lista de símbolos
% \begin{simbolos}
%   \item[$ \Gamma $] Letra grega Gama
%   \item[$ \Lambda $] Lambda
%   \item[$ \zeta $] Letra grega minúscula zeta
%   \item[$ \in $] Pertence
% \end{simbolos}

% Inserir o SUMÁRIO
\pdfbookmark[0]{\contentsname}{toc}
\tableofcontents*
\cleardoublepage

% ----------------------------------------------------------
% ELEMENTOS TEXTUAIS (Capítulos)
% ----------------------------------------------------------
\textual
% Elementos textuais com numeração arábica
\pagenumbering{arabic}
% Reinicia a contagem do número de páginas
\setcounter{page}{1}

\sethlcolor{Blue}

% Inclui cada capitulo 
